\newpage
\subsection{Limitations}
\label{sec:limitations}

The results must be considered within the limits of the experiment.

The lab is small and single-vendor: four Nokia SR Linux nodes and four clients, with exactly
one injected fault at a time. An enterprise network has many more devices,
mixed vendors, concurrent faults and noisy alarms. Nothing here shows how the system scales
to that. The results are a feasibility statement for small environments and not for big enterprise networks.
The runs are also non-interactive, each consists of a single symptom
description with no follow-up dialogue (\Cref{sec:experiment-execution}). In a real use case,
a network operator could answer questions the agents raise, add observations, or steer the
investigation. Such help would plausibly raise the success rate, particularly for the models
that stopped at a plausible but wrong explanation. The reported numbers measure the fully
autonomous mode.

Three restrictions have to be considered regarding the measurements. The time and cost savings compare against an estimate
of how long an engineer would need. No human trials were run. Each scenario--model pair was
repeated only five times, so one additional or one fewer successful run shifts a rate by 20 percentage points.
And the decision whether a run was a success or a miss was made by one person against a fixed criterion. Borderline answers
could be judged differently by another person.

Two further restrictions concern the recorded data. First, GLM-5.2's success rate may be
understated. In 4 of its 50 runs the model ended the agent loop after a single response
without calling any tools (the reasoning stated that tools should be called, but none were), so no investigation took place (\Cref{sec:results-correctness}),
all four count as misses. Excluding these degenerate runs, its found rate would be 37/46
($\approx$80\,\%) instead of 74\,\%. GLM-5.2 had been released not long before the
experiment window, so issues in its tool-calling layer are a plausible cause. A follow-up
check on a self-hosted GLM-5.2 instance at FHNW did not reproduce the behaviour, which
points towards the serving layer of the API provider rather than the model itself. But this
check covered only a very small set of runs. Since the failure appeared in 4 of 50 runs
(8\,\%) in the main experiment, its absence in such a small sample could also be
chance, so this observation is an indication and not proof. This means the reported 74\,\% may understate the model's actual diagnostic ability.
Second, runs were capped at 600 seconds, and this cap does not weigh equally on all models.
The cap also shapes the difficulty answer for the mid tier: 16 of the mid-tier models' 22
misses are cap hits rather than wrong conclusions, so with a longer time budget their picture
on the medium and hard scenarios could look better.
Token generation speed differs between models, and for the open-weight models it
is partly a property of the serving provider (compute power) rather than of the model itself. A slowly served
model reaches the time limit, without being able to do as much work as a fast one, independent
of its diagnostic ability. DeepSeek-V3.2 is the clearest case. A large part of its timeouts can be attributed to slow
token generation rather than to a deeper investigation, since its tool-call count is close to
GLM-5.2's (\Cref{fig:tool-calls}) while far more of its runs hit the cap, so its low
found rate partly reflects serving speed instead of only diagnostic capability. Its true time
and success numbers beyond the cap are unknown (more than half of its runs were cut off). Token and cost figures for DNF
runs are even a bit undercounted. Usage is recorded only when an agent run returns, so the cap
discards the orchestrator's own consumption and that of any sub-agent still working,
counting only the sub-agents that had already completed. The reported cost of models with many
DNFs understates what they actually consumed.
